%------------------------------------------------------------------------------%
\begin{exercises}

%------------------------------------------------------------------------------%
\begin{exercise}
Use a parte 1 do Teorema Fundamental do Cálculo para calcular a
derivada das seguintes funções.
\begin{mathlist}
  \item g(x)=\int_1^x \frac{1}{1+t^3} dt
  \item h(x)=\int_x^1 \cos\left(\sqrt{t}\right)dt
  \item h(x)=\int_1^{e^x} \ln(t)dt
  \item p(x)=\int_1^{\sqrt{x}} \frac{z^2}{z^4+1} dz
  \item j(x)=\int_{\sin(x)}^1  \sqrt{1+v^2}dv
  \item h(x)=\int_{1-3x}^{1} \frac{u^3}{1+u^2} du
  \item h(x)=\int_{2x}^{3x} \frac{u^2-1}{1+u^2} du
  \item h(x)=\int_{x}^{x^2} e^{t^2} dt
  \item h(x)=\int_{\sqrt{x}}^{2x} \arctan(t)dt
  \item p(x)=\int_{\sin(x)}^{\cos(x)} ln(1+2v) dv
\end{mathlist}
\begin{answer}
  \begin{mathlist}
    \item g'(x)=\frac{1}{1+x^3}
    \item h'(x)=-\cos(\sqrt{x})
    \item \frac{x^3}{3}-\frac{1}{x}+c
    \item p'(x)=\frac{1}{2\sqrt{x}}\left( \frac{x^2}{x^4+1} \right)+ \frac{9}{2}u^2+4u+C
    \item \frac{1}{6}v^6 +v^4+2v^2+C
    \item h'(x)=-3 \left( \frac{(1-3x)^3}{1+(1-3x)^2}\right)
    \item \frac{-8x^2+2}{1+4x^2}+\frac{27x^2-3}{1+9x^2}
    \item -e^{x^2}+2xe^{x^4}
    \item -\frac{1}{2\sqrt{x}}\left( \arctan(\sqrt{x}) \right)+2 \arctan(2x)
    \item -\cos(x)ln(1+2\sin(x))-\sin(x) ln(1+2\cos(x))
  \end{mathlist}
\end{answer}
\end{exercise}
%------------------------------------------------------------------------------%

%------------------------------------------------------------------------------%
\begin{exercise}
Seja \(f(x)=\int_0^x (1-t^2)e^{t^2} dt\)
\begin{alphalist}
  \item Encontre  e classifique os pontos críticos de $f$.
  \item Analise o crescimento e o decrescimento de $f$.
\end{alphalist}
\begin{answer}
  \begin{alphalist}
    \item Pontos críticos: $x=-1$ e $x=1$.
    \item $f$ é crescente no intervalo $(-1,1)$ e decrescente nos
          intervalos $(-\infty,-1)$ e $(1, \infty)$.
  \end{alphalist}
\end{answer}
\end{exercise}
%------------------------------------------------------------------------------%

%------------------------------------------------------------------------------%
\begin{exercise}
Em qual intervalo a curva \(y= \int_0^x \frac{t^2}{t^2+t+2} dt\)
é côncava para baixo?
\begin{answer}
  No intervalo $(-4,0)$
\end{answer}
\end{exercise}
%------------------------------------------------------------------------------%

%------------------------------------------------------------------------------%
\begin{exercise}
Calcule $f'(2)$ se $f(x)= e^{g(x)}$ e \(g(x)= \int_2^x \frac{t}{1+t^4}dt\).
\begin{answer}
  $f'(2)=\frac{2}{17}$
\end{answer}
\end{exercise}
%------------------------------------------------------------------------------%

%------------------------------------------------------------------------------%
\begin{exercise}
Se \(f(x)= \int_0^{\sin(t)} \sqrt{1+t^2} dt\)
e  \(g(y)= \int_3^y f(x)dx\),
encontre \(g''\left(\sfrac{\pi}{6}\right)\).
\begin{answer}
  $\frac{\sqrt{15}}{4}$
\end{answer}
\end{exercise}
%------------------------------------------------------------------------------%

%------------------------------------------------------------------------------%
\begin{exercise}
Calcule a integral definida
\begin{mathlist}
  \item \int_1^4 5-2t+3t^2\,dt
  \item \int_1^9 \sqrt{x}\,dx
  \item \int_0^1(4-x) \sqrt{x}\,dx
  \item \int_1^8 \frac{1}{x^{\sfrac{2}{3}}}\,dx
  \item \int_{\sfrac{\pi}{6}}^{\pi}\sin(x) dx
  \item \int_{0}^{\frac{\pi}{4}}  \sec(v)\tan(v) dv
  \item \int_{1}^{2} \frac{4+u^2}{u^3} du
  \item \int_{1}^{3} 2\sin(x)-e^x dx
  \item \int_{\sfrac{1}{2}}^{\sfrac{1}{\sqrt{2}}} \frac{4}{\sqrt{1-x^2}} dx
  \item \int_{-1}^{1} e^{t+1}dt
  \item \int_{\sfrac{1}{\sqrt{3}}}^{\sqrt{3}} \frac{8}{1+v^2} dv
  \item \int_0^{\pi} f(x)dx \)
        se \(
              f(x) = \begin{cases}
                       \sin(x) & 0 \leq x < \sfrac{\pi}{2} \\
                       \cos(x) & \sfrac{\pi}{2}  \leq x  \leq \pi
                     \end{cases}
           \)
  \item \int_{-2}^{2} f(x)dx \)
        se \(
              f(x) = \begin{cases}
                        2     &     -2 \leq x \leq 0 \\
                        4-x^2 & \hpm 0 \leq x \leq 2
                     \end{cases}
           \)
  \item \int_0^1 5x-5^x dx
  \item \int_0^1 x\big(\sqrt[3]{x}+ \sqrt[4]{x}\big)dx
  \item \int_0^{\sfrac{\pi}{3}} \frac{\sin(\theta)+\sin(\theta) \tan(\theta)}{\sec^2(\theta)} d\theta
  \item \int_0^{\sfrac{1}{\sqrt{3}}} \frac{t^2-1}{t^4-1}dt
  \item \int_1^4 \frac{\sqrt{y}-y}{y^2}dy
  \item \int_0^2 \abs{2x-1} dx
  \item \int_1^2 \frac{1}{x^2}-\frac{4}{x^3}\,dx
\end{mathlist}
\begin{answer}
  \begin{mathlist}
    \item 63
    \item \frac{52}{3}
    \item \frac{128}{15}
    \item 3
    \item 1+ \frac{\sqrt{3}}{2}
    \item \frac{3}{2}+ \ln(2)
    \item 3-2\cos(3)+e^3
    \item \frac{\pi}{3}
    \item e^2-1
    \item \frac{4\pi}{3}
    \item \frac{28}{3}
    \item \frac{5}{2}-\frac{4}{\ln(5)}
    \item \frac{55}{63}
    \item \frac{1}{2}
    \item \frac{\pi}{6}
    \item 1-\ln(4)
    \item -\frac{7}{2}
    \item -1
  \end{mathlist}
\end{answer}
\end{exercise}
%------------------------------------------------------------------------------%

%------------------------------------------------------------------------------%
\begin{exercise}
Determine o que está errado em cada na equação.
\begin{mathlist}
  \item
    \int_{-2}^1 x^{-4} dx
    = \frac{x^{-3}}{-3} \evalat{-2}{1}
    = - \frac{-3}{8}
  \item
    \int_{0}^{\pi} \sec^2(x) dx
    = \tan(x) \evalat{0}{\pi}
    = 0
\end{mathlist}
\end{exercise}
%------------------------------------------------------------------------------%

%------------------------------------------------------------------------------%
\begin{exercise}
Calcule as integrais a seguir, interpretando-a em termos das áreas,
faça um esboço da região.
\begin{mathlist}
  \item \int_{-1}^2 1-x\,dx
  \item \int_{-3}^0 1+\sqrt{9-x^2}\,dx
\end{mathlist}
\end{exercise}
%------------------------------------------------------------------------------%

%------------------------------------------------------------------------------%
\begin{exercise}
Nos exercícios a seguir, encontre:
\begin{romanlist}
    \item O valor médio de $f$ no intervalo dado.
    \item O valor de $c$ tal que $f_{\text{médio}}=f(c)$.
    \item Esboce o gráfico de $f$ e um retângulo cuja área é a
          mesma que a área sob o gráfico de $f$ nesse intervalo.
\end{romanlist}%
\begin{mathlist}
  \item f(x)=(x-3)^2,    \) \tabto{50mm} \([2, 5]\)
  \item f(x)=\sqrt{x},   \) \tabto{50mm} \([0, 4]\)
  \item f(x)=-3x^2-1,    \) \tabto{50mm} \([0, 1]\)
  \item f(x)= \abs{x}-1, \) \tabto{50mm} \([-1, 1]\)
\end{mathlist}
\begin{answer}
\begin{alphalist}
  \item
    \begin{romanlist}[2]
      \item \(1\)
      \item \(c = 2\) e \(c = 4\)
    \end{romanlist}
  \item
    \begin{romanlist}[2]
      \item \(\frac{4}{3}\)
      \item \(c=\frac{16}{9}\) e \(c=4\)
    \end{romanlist}
  \item
    \begin{romanlist}[2]
      \item \(-2\)
      \item \(c=2\)
    \end{romanlist}
  \item
    \begin{romanlist}[2]
      \item \(\frac{-1}{2}\)
      \item \(c=-\frac{1}{2}\) e \(c=\frac{1}{2}\)
    \end{romanlist}
\end{alphalist}
\end{answer}
\end{exercise}
%------------------------------------------------------------------------------%

%------------------------------------------------------------------------------%
\begin{exercise}
Use a propriedade 10 de integrais para estimar o valor da integral.
\begin{mathlist}
  \item \int_{0}^2 xe^{-x} dx
  \item \int_{\sfrac{\pi}{4}}^{\sfrac{\pi}{3}} \tan(x)dx
\end{mathlist}
\begin{answer}
  \begin{mathlist}
    \item 0
          \leq \int_{0}^2 xe^{-x} dx
          \leq \frac{1}{e}
    \item \frac{\pi}{12}
          \leq \int_{\sfrac{\pi}{4}}^{\sfrac{\pi}{3}} \tan(x)dx
          \leq \frac{\pi}{2} \sqrt{3}
  \end{mathlist}
\end{answer}
\end{exercise}
%------------------------------------------------------------------------------%

%------------------------------------------------------------------------------%
%\begin{exercise}
%\begin{answer}
%\end{answer}
%\end{exercise}

%------------------------------------------------------------------------------%
%\begin{exercise}
%\begin{mathlist}[2]
%  \item
%  \item
%  \item
%\end{mathlist}
%\begin{answer}
%  \begin{alphalist}[3]
  %    \item
  %    \item
  %    \item
  %  \end{alphalist}
%\end{answer}
%\end{exercise}

\end{exercises}
%------------------------------------------------------------------------------%
